\subsection{WAGON COUNTERS}

Wagon counters are used to simulate the "soft" logistical "tail" found in all armies. They are slow, vulnerable, and tie up useful in their protection. But they are absolutely essential to any army's effectiveness and survival.

\begin{enumerate}
  \item Wagon counters have a Movement Allowance Factor of "4".

  \item One Wagon counter can "carry" one Supply counter.

  \item In Combat, Wagon counters have no value. If there are no friendly Combat units (Infantry, Cavalry, Horse Artillery or Partisan units) in a hex with a Wagon counter, and it is placed in a Phasing enemy unit's Zone of Control, the Wagon counter is considered to be captured, and can immediatley be replaced by one of the enemy's Wagon counters (to show that it has switched ownership). This alos applies if Combat units originally in the hex with the Wagon counter have all been eliminated. As soon as a Wagon counter has been captured, it may be immediately moved up to its full Movement Allowance Factor by the captor.

  \item Wagon counters can be destroyed by Phasing Infantry and/or Cavalry units that are in the same hex.

  \item Combat units in the same hex as the Wagon counter do not have the option to withdraw. They must fight at least one Round of Combat before withdrawing. If forced to retreat, the wagons can retreat with the Combat units, after the one Round of Combat.
\end{enumerate}

\subsection{SUPPLY COUNTERS}

Supply counters in SHENANDOAH are used largely to simulate ammunition supplies, although to a lesser extent they also simulate supplies of other essential items such as equipment, medicine, rations, etc, these being considered to be items that the troops could "forage" for if the need arose.

\begin{enumerate}
  \item Supply counters have no Movement Allowance Factor. They can only be moved along operable rail lines, or by being "carried" by a Wagon counter.

  \item Other units may be moved into a hex containing a Supply counter without the need to expend Movement Allowance Factors to Stack or Unstack with it. However, it is necessary for a Wagon counter that moves into a hex with a Supply counter to expend Movement Allowance Factors to Stack with it before the Wagon counter can "carry" the Supply counter. The reverse is true for the Wagon counter to move on and leave the Supply counter behind - Movement Allowance Factors must be expended for Unstacking.

  \item To be "supplied", a Combat unit must be in the same hex as a friendly Supply counter, or able to trace a path to one that is four hexes or less away. A path traced to a friendly Supply counter cannot run through a "ridge" hexside, through a river hexside during a Turn of Mud (except over an undestroyed bridge), or through a hex that contains an enemy unit. The path can be traced through other types of terrain (including through enemy Zones of Contro). Units that cannot meet the above conditions are considered to be "unsupplied". Supply considerations do not apply to Wagon counters.

  \item Supplied units may move and have Combat normally. Unsupplied units are at a disadvantage in Combat (see the "Unsupplied" column of the COMBAT TABLE MODIFIERS CHART). Unsupplied units also lose two from their regular Movement Allowance Factor - this applies if they are over four hexes from a friendly Supply counter when they start their movement; you are allowed to move up a Supply counter so as to supply these units before you begin moving them. If any unit in an attack is unsupplied, all units in that attack suffer the out of supply modifiers.

  \item Supply counters can supply any number of Combat units in any number of hexes, provided that all those units can trace a "safe" four-hex route back to the Supply counter.

  \item When the Phasing Player attacks with supplied Combat units, their Supply counter is removed from the Mapboard (it has been used to "supply" the attack - it has not been destroyed, and the opposing player will get no Victory Points for its removal). One supply counter may supply any number of attacks within its four hex radius, and any number of Rounds of Combat that may take place during that Phase of the Turn. The Non-Phasing Player does not remove any Supply counters, even if he counter-attacks.

  \item Note that the unsupplied units of the Non-Phasing Player suffer disadvantages in Combat, as well as the Phasing Player's units. Also, note that it is possible for the Player who goes first during a Turn to make supplied attacks during his own Phase, then remove his Supply counter and, unless he has another Supply counter within four hexes, to be unsupplied during the other Player's phase.

  \item Supply counters can be destroyed by Infantry and/or Cavalry units that are in the same hex.

  \item In Combat, Supply counters have no value. If there are no friendly Combat units in their hex, or if all friendly Combat units in the hex have been eliminated, and the Supply counter is placed in a Phasing enemy unit's Zone of Control, it is considered captured, and can be immediately replaced by one of the enemy's Supply counters.

  \item The Phasing Player must use a Supply counter when attacking if one is close enough to be used - he cannot choose to attack "unsupplied" in order to keep the Supply counter for later Turns. Also, the Phasing Player is not allowed to "use" more Supply counters than necessary to supply his attacks - i.e. there can never be a case where two Supply units so placed as to supply the same attacks during a Turn can both be "used" to supply them. A player must destroy excess supply counters, he cannot "use" more than are necessary.

  \item With the exception of Partisan units (who are merely considered to be "unsupplied"), all Combat units that are "unsupplied" while in a Devastated hex will one one Combat Factor of deserters per Stack per Turn this situation exists. These deserters will count for Victory Points, just as if they had been eliminated in Combat.

  \item Artillery units that are unsupplied have no combat value, i.e. their Combat Factor is "zero".

  \item If the Phasing Player enters an enemy Zone of Control, but makes no attacks during his segment of the Rounds of Combat, no supply unit is used.
\end{enumerate}

\subsubsection{ENTERING \& EXITING WAGON AND SUPPLY COUNTERS}

\begin{enumerate}
  \item New Supply and/or Wagon counters can enter the Mapboard at any time the player wishes to bring them in. They are placed in the appropriate box of the TURN RECORD CHART - they must be placed there for the Turn following the current one.

  \item New Supply and/or Wagon counters can only enter the Mapboard at designated entry hexes, or within four hexes of them along the same edge of the Mapboard.

  \item Introducing new Supply and/or Wagon counters gives the opposing player Victory Points.
\end{enumerate}